% ---------------------------------------------------------------------------
% A clean document: title, abstract, Figure 1. Nothing else.
%
% main.tex carries the full record --- five results sections, the campaign
% deck, 15k words of experimental notebook --- and that is the right home for
% it. This file exists to build the paper forward one figure at a time without
% the rest of the draft in the way. Each figure gets added here when it is
% settled; nothing is deleted from main.tex to make room.
%
% \draftmodefalse, so no \plan block or \todo marker can reach this PDF. That
% is deliberate: an earlier abstract shipped bracketed placeholders into the
% rendered output for weeks because draft mode was left on.
%
% Build:  make clean-pdf  ->  clean.pdf
% ---------------------------------------------------------------------------
\documentclass[11pt]{article}

% --- Packages ---
\usepackage[utf8]{inputenc}
\usepackage[T1]{fontenc}
\usepackage{lmodern}
\usepackage[margin=1in]{geometry}
\usepackage{microtype}

\usepackage{graphicx}
\usepackage{xcolor}
\usepackage{amsmath, amssymb}
\usepackage{siunitx}
% siunitx drops a leading "+" unless told otherwise, so \SI{+23.0}{points} set
% "23.0 points" -- turning a signed margin into a level and silently reversing
% the reading of any negative-looking cell beside it. There is no unit called
% "points" either; declare it rather than letting siunitx guess.
\sisetup{retain-explicit-plus}
\DeclareSIUnit{\points}{points}
\usepackage{booktabs}
\usepackage{caption}
\usepackage[numbers,sort&compress]{natbib}
\usepackage{hyperref}
\hypersetup{colorlinks=true, linkcolor=black, citecolor=black, urlcolor=blue}

% authblk is loaded only if installed. Nothing here uses \affil -- it was
% pulled in for author affiliations that were never written -- and on a
% machine whose texlive lacks it the whole document stopped building over
% a package it does not use.
\IfFileExists{authblk.sty}{\usepackage{authblk}}{}

% Allow hyphenation inside \texttt{}. Without this, snake_case identifiers
% (mutation_diversity, consensus_divergence, ...) are unbreakable and push
% lines into the margin. [htt] enables breaking in typewriter text.
% Also optional: it only permits hyphenation inside \texttt, so a texlive
% without it loses slightly worse line breaks around identifiers and
% nothing else.
\IfFileExists{hyphenat.sty}{\usepackage[htt]{hyphenat}}{}

\graphicspath{{figures/pdf/}}

% --- Shared color palette ---
% Author's own "good"/IPD family -- see paper/figures/palette.py for the
% order they must be used in and which pairs conflict.
\definecolor{cleoBlue}{HTML}{4499E7}
\definecolor{cleoYellow}{HTML}{FAC72C}
\definecolor{cleoPurple}{HTML}{B41F85}
\definecolor{cleoRed}{HTML}{E44A3E}
\definecolor{cleoOrange}{HTML}{FF8C00}
\definecolor{cleoTeal}{HTML}{28B0C1}
\definecolor{cleoGreen}{HTML}{AAC32F}
\definecolor{cleoGray}{HTML}{6B7280}

% --- Convenience macros ---
\newcommand{\cleo}{\textsc{cleo}}
\newcommand{\todo}[1]{\textcolor{cleoRed}{\textbf{[TODO: #1]}}}

% --- Draft / outline mode ---
% While the paper is a scaffold, the planning content is the document. Set
% \draftmodefalse to hide every \plan block and \todo marker in one shot,
% leaving only real prose — do that when preparing a submission draft.
\newif\ifdraftmode
\draftmodetrue

% Internal links (table of contents, cross-references) render in blue while
% drafting so the contents page reads as clickable. linktoc=all makes both the
% entry text and its page number active, not just the number. Submission mode
% keeps them black — see the \hypersetup on the hyperref line above.
\ifdraftmode
  \hypersetup{linktoc=all, linkcolor=cleoBlue}
\fi

% \plan{...} renders a planning note. In draft mode it is visible and set in
% gray; otherwise it vanishes.
\newcommand{\plan}[1]{%
  \ifdraftmode
    \par\vspace{0.3em}%
    \begingroup\small\color{cleoGray}\setlength{\leftskip}{1.2em}%
    #1\par\endgroup\vspace{0.3em}%
  \fi}

% planlist: bulleted planning notes. Same visibility rules as \plan.
% `comment` is only needed by the \draftmodefalse branch below, and it is not
% in every texlive. Where it is missing, planlist bodies are typeset into a
% box that is then dropped -- the content does not reach the page either way.
\IfFileExists{comment.sty}{\usepackage{comment}\newif\ifhavecomment
  \havecommenttrue}{\newif\ifhavecomment\havecommentfalse}
\newenvironment{planlistbody}%
  {\begingroup\small\color{cleoGray}%
   \begin{itemize}\setlength{\itemsep}{0pt}\setlength{\parskip}{0pt}}%
  {\end{itemize}\endgroup}
\ifdraftmode
  \newenvironment{planlist}{\begin{planlistbody}}{\end{planlistbody}}
\else
  \ifhavecomment
    \excludecomment{planlist}
  \else
    \newenvironment{planlist}{\setbox0\vbox\bgroup}{\egroup}
  \fi
\fi

% --- Status tags for the document index ---
\newcommand{\statusdone}{\textcolor{cleoGreen}{\textbf{drafted}}}
\newcommand{\statusscaffold}{\textcolor{cleoOrange}{\textbf{scaffold}}}
\newcommand{\statusblocked}{\textcolor{cleoRed}{\textbf{needs data}}}

% Placeholder for an experiment result not yet measured. Visually distinct from
% \todo so the run-log table reads at a glance.
\newcommand{\todoresult}{\textcolor{cleoGray}{\textit{pending}}}

\draftmodefalse

\title{Design of high fidelity combinatorial libraries}
\author{}
\date{\today}

\begin{document}
\maketitle

\begin{abstract}
\noindent
Recent tools for de novo design have revolutionized the way we propose novel
proteins and the types of challenges we are able to tackle. While these tools
have dramatically increased the hit rate of functional designs, zero-shot
generations often require further optimization to enhance the desired
properties. We present a novel method for generating high fidelity diverse
combinatorial libraries. Traditional library generation methods rely on error
prone PCR or natural fragment shuffling to generate diversity. Such approaches
do not take into account the filtering of designs through structure prediction
models or other desired oracles. We demonstrate the ability to post-train
inverse folding models to design fragments with filter pass rates of 93.4\% and
libraries which can easily scale beyond $10^{12}$ individually addressable
variants. Our post-training framework enables control over the search space by
centering the library around a reference sequence with tunable mutation
distance. Once designed, we demonstrate how individual variants are assembled
with simple molecular biology techniques (golden gate assembly). Our cell-free
assay platform enables us to screen enzymes at the scale of $10^{3}$
measurements per day, each with replicates. The libraries generated are easily
much larger than most traditional and modern high throughput assays are able to
screen, so we demonstrate how the library can be sub sampled smartly to
optimize enzymes. For a de novo designed PETase we demonstrate our platform
finds a variant with 250x increase in catalytic activity.
\end{abstract}

% Figure 1's legend sits on the page AFTER the figure.
%
% The float carries no \caption, so it does not step the figure counter; the
% \captionof below does, and the numbering comes out the same. Done this way
% rather than with an empty second float, which WOULD step the counter and
% silently renumber the method figure.
\begin{figure}[p]
  \centering
  \includegraphics[width=\textwidth]{concept_library.pdf}
\end{figure}

\clearpage
\noindent\captionof{figure}{\textbf{A combinatorial library, and the loop around it.}
  \textcolor{cleoRed}{[PLACEHOLDER --- the prose is yours. Kept short so the
  float fits; the facts below are the ones a reader cannot get from the
  picture.]}
  \textbf{(a)}~One toy landscape searched three ways with the same number of
  variants built and screened ($N=36$). Illustrative: no axis is measured.
  \textbf{(b)}~eBlocks are billed at a 300-bp minimum, which binds from $k=2$
  on, so past that point cutting finer stops buying shorter DNA and only buys
  more pieces of it. $10^{12}$ variants costs \$42T as whole genes, \$42M at
  $k=2$, \$84k at $k=4$, \$5.3k at $k=8$.
  \textbf{(c)}~\textcolor{cleoRed}{[Reserved --- to be filled in.]}
  \textbf{(d)}~What makes the library combinatorial. One fold is cut at $k-1$
  boundaries; each slot is then an independent pool of $P$ parts; and any one
  part per slot assembles a full-length design. Ordering $k\times P$ fragments
  therefore builds $P^{k}$ sequences, which is the exponent panel~(b) prices.
  Parts per slot are drawn as three to keep the columns legible.
  \textbf{(e)}~The experimental loop. The arrow returns to the \emph{selector},
  not the library: parts are synthesised once and each round changes only
  which of them get combined.}
\label{fig:concept}
\clearpage

\noindent\textcolor{cleoRed}{[PLACEHOLDER --- the prose is yours. The
numbers below are from \texttt{paper/figures/panel\_cost.py}, which prints
them; do not retype them by hand.]}

Post-training costs a fixed number of folds. A run is 200 steps, and
\texttt{build\_chimeras} emits $\text{rotations}\times\text{batch}=64$
chimeras per step, each folded best-of-5: $4\times16\times5=320$ folds per
step, \SI{64000}{} in total. The bill does not grow with $k$ --- cutting a
design finer does not fold more of it --- so one figure covers the whole
sweep. The alternative is to fold your way to a usable library: generate with
stock ProteinMPNN, fold every candidate, and keep what passes. At vanilla
pass rate $p$, obtaining $S$ members that fold costs $5S/p$ folds, which is
linear in $S$, where post-training is a constant \SI{64000}{} \emph{however
large the library is} --- the enrichment lives in the policy, not in the
members, so anything drawn from a trained policy inherits the higher pass
rate without being folded at all.

Below a few thousand members brute force is the cheaper option, and the
break-even depends on how good the untrained model already was: \SI{3300}{}
members at M0255's \SI{25.8}{\percent}, \SI{6300}{} at M0315's
\SI{49.6}{\percent}, \SI{7700}{} at M0097's \SI{59.8}{\percent}. That
range is the whole point. A combinatorial library is not a few thousand
members: $k=4$ with 100 parts per slot is $10^{8}$, where filtering by
folding would cost ${\sim}8{\times}10^{8}$ folds --- four orders of
magnitude more than training once --- and at $10^{12}$ the comparison stops
being a comparison.

% Figure 2 is the computational result, and its legend follows it the way
% Figure 1's does.
\begin{figure}[p]
  \centering
  \includegraphics[width=\textwidth]{fig2_libraries.pdf}
\end{figure}

\clearpage
\noindent\captionof{figure}{\textbf{Designing high-fidelity combinatorial
  libraries.} \textcolor{cleoRed}{[PLACEHOLDER --- the prose is yours.]}
  \textbf{(a)}~How post-training works here. ProteinMPNN is pretrained and
  emits whole sequences, so nothing in it says which fragments
  \emph{recombine}; the method manufactures that signal. A batch is cut into
  $k$ slots, chimeras are built by swapping slots between parents, each
  chimera is folded, and every fragment is given the mean outcome of the
  chimeras it appeared in --- here one fragment occurs in three, two of which
  fold, so its reward is $2/3$. That per-fragment reward, not the parent's
  score, drives the GRPO update. Folding during training is best-of-5 against
  a \emph{pool} of RF3, Boltz and Protenix under a round robin; AlphaFold3 is
  held out of training entirely and scores every number in (b)--(d).
  One row per backbone throughout.
  \textbf{(b)}~One row per backbone, read left to right: the target, what
  post-training buys it, and how far that holds as the library is cut finer.
  \emph{Left:} the three targets, coloured by their four fragments, with the
  motif RFdiffusion2 holds fixed drawn as sticks and the ligand in grey. The
  cut lines run \emph{through} the active site: under the same equal-count
  split, M0097's motif falls on three different fragments (A26, A59/A61,
  A93) and M0255's and M0315's each fall on two. Every recombination below is
  therefore reassembling an active site from parts contributed by different
  designs.
  \emph{Middle:} what post-training buys, measured on recombinants --- decoded
  from MPNN, cut at $k=4$, shuffled across parents, reassembled, folded,
  never on parents. Arrows run from vanilla to post-trained. Success rises on
  all three, by $+26.2$, $+60.2$ and $+48.8$ points. The vanilla arm is not a
  single temperature: its fragments come from parents sampled at
  $T=0.1$--$1.0$ and then filtered to those that pass, which skews the pool
  toward low $T$, so it is the filtered union of the whole sweep rather than
  any one setting. Distinct mutations are counted over \emph{passing}
  recombinants only and rarefied to a matched survivor count, so a wider
  spread is genuinely wider coverage and not just a larger sample; two
  backbones move right as well as up, and M0315 trades 193 mutations for its
  48.8 points.
  \emph{Right:} a policy
  trained at each $k$ and recombined at that same $k$, with every cell run
  for a matched 200 steps so training budget cannot covary with $k$.
  Retention is flat through $k=16$ --- 11-residue fragments --- and falls
  away below that. The control --- one $k=4$ policy re-cut at every $k$, with no
  training at the finer granularity --- is in Fig.~\ref{fig:posthoc}.
  \textcolor{cleoRed}{Caveats that must survive into the text: every cell is
  single-seeded against a documented 21-point run-to-run swing; M0255 at
  $k=36$ and $k=64$ was still climbing at step 200, so those are lower
  bounds.}}
\label{fig:libraries}
\clearpage

\begin{figure}[p]
  \centering
  \includegraphics[width=\textwidth]{fig3_momi.pdf}
\end{figure}

\clearpage
\noindent\captionof{figure}{\textbf{A combinatorial PETase library, built
  and run.} \textcolor{cleoRed}{[PLACEHOLDER --- the prose is yours.]}
  The second of two libraries on the same G20 parent; the two share 5 parts
  out of 3,629, so they are separate designs rather than one campaign
  continued.
  \textbf{(a)}~The 209-aa G20 parent cut into $k=5$ slots, coloured by slot,
  with the positions no design ever mutates in grey. Slot windows are
  measured rather than assumed: the parent is the positive control
  (\texttt{1.g20\_\_\_\dots\_\_\_5.g20}, 52 replicate wells), and the
  designs that swap exactly one slot away from it show which positions move
  --- residues 1--40, 45--84, 89--126, 131--163 and 168--209, non-overlapping
  and in order. \emph{Right:} the catalytic triad is \emph{split} across
  slots, Ser66 on slot 2 and His107/Asp122 on slot 3, with the PET substrate
  in grey and the two catalytic hydrogen bonds dashed. Every recombination in
  this library therefore reassembles the triad across a cut.
  \textbf{(b)}~How far each slot's parts sit from the G20 part they replace,
  one distribution per slot over the $\sim$700 distinct parts it carries.
  Slot~1 roams furthest (median 9, up to 17 mutations); slot~4 least (median
  6, never more than 9).
  \textbf{(c)}~Activity by round, as a multiple of the G20 parent measured on
  the same plates. The median goes $0.83 \rightarrow 0.36 \rightarrow 5.54
  \rightarrow 8.06$, and the share of designs beating G20 goes
  \SI{38.7}{\percent} $\rightarrow$ \SI{32.1}{\percent} $\rightarrow$
  \SI{92.9}{\percent} $\rightarrow$ \SI{99.7}{\percent} --- 368 of 369
  designs in the final round. The best single design is
  $18.98\times$ G20.
  \textcolor{cleoRed}{Caveats: round~1 is a dip, not noise, and the text
  should say why; and activity here is the plate screen, which is not the
  same measurement as kinetics.}}
\label{fig:momi}
\clearpage

% --- supplementary -------------------------------------------------------
% Numbered S1, S2, ... by resetting the counter and reformatting it, so the
% main figures keep 1 and 2.
\setcounter{figure}{0}
\renewcommand{\thefigure}{S\arabic{figure}}

\noindent{\large\textbf{Supplementary figures}}

\begin{figure}[p]
  \centering
  \includegraphics[width=\textwidth]{figS_posthoc.pdf}
\end{figure}

\clearpage
\noindent\captionof{figure}{\textbf{Training at $k$, against the control
  dropped from Fig.~\ref{fig:libraries}.}
  \textcolor{cleoRed}{[PLACEHOLDER --- the prose is yours.]}
  \textbf{(a)}~Retention against the number of fragments per design, one
  facet per backbone. Blue: a policy trained at that $k$ and recombined at
  the same $k$. Grey: a single policy trained at $k=4$ and re-cut at every
  $k$, with no training at the finer granularity. Every \emph{trained} cell
  ran a matched 200 steps, so training length cannot covary with $k$; the
  control needs no training at all, which is what makes it the right
  counterfactual.
  \textbf{(b)}~The difference between the two, which widens monotonically
  with $k$: $+8.3$ points at $k=2$, $+18.0$ at $k=8$, $+27.9$ at $k=16$ and
  $+32.2$ at $k=36$. Cutting a finished design finer loses recombinants
  faster than training a policy to be cut that finely does, and the finer the
  cut the more training at that granularity is worth.
  \textcolor{cleoRed}{Caveats: every cell is single-seeded against a
  documented 21-point run-to-run swing; M0255 at $k=36$ and $k=64$ was still
  climbing at step 200, so those are lower bounds.}}
\label{fig:posthoc}
\clearpage

\end{document}
